% ch4_c (书页 175-186 / p190-p201)
%--B-TAIL--
%JOIN%
$T_{\mu\nu}t^\mu t^\nu$ 就对所有类时矢量 $t^\mu$ 非负（请自己验证这一点；这不过是矢量相加而已）。因此我们来计算
\begin{equation*}
T_{\mu\nu}U^\mu U^\nu = \rho, \qquad T_{\mu\nu}\ell^\mu\ell^\nu = (\rho + p)(U_\mu \ell^\mu)^2.
\tag{4.108}
\end{equation*}
因此 WEC 意味着 $\rho \geq 0$ 和 $\rho + p \geq 0$。这些听起来都十分合理：能量密度非负，并且压强相对于能量密度不能太大。当然，我们不必把自己限制在理想流体上，我们只是利用它们来洞察能量条件所施加的要求。

能量条件有许多种，各自适用于不同的情形。下面是一些最常见的：

\begin{itemize}
\item \textbf{弱能量条件}（Weak Energy Condition），或称 WEC，正如刚才讨论的，它要求对所有类时矢量 $t^\mu$ 有 $T_{\mu\nu}t^\mu t^\nu \geq 0$，或者等价地，$\rho \geq 0$ 且 $\rho + p \geq 0$。

\item \textbf{类光能量条件}（Null Energy Condition），或称 NEC，要求对所有类光矢量 $\ell^\mu$ 有 $T_{\mu\nu}\ell^\mu\ell^\nu \geq 0$，或者等价地，$\rho + p \geq 0$。它是 WEC 的特例，只是把类时矢量换成了类光矢量。此时能量密度可以取负值，只要存在补偿性的正压强。

\item \textbf{主能量条件}（Dominant Energy Condition），或称 DEC，包含 WEC（对所有类时矢量 $t^\mu$ 有 $T_{\mu\nu}t^\mu t^\nu \geq 0$），外加一条附加要求：$T^\mu{}_\nu t^\nu$ 是非类空（nonspacelike）矢量（也就是说，$T_{\mu\nu}T^\nu{}_\lambda t^\mu t^\lambda \leq 0$）。对于理想流体，这些条件合在一起等价于一个简单的要求：$\rho \geq |p|$；能量密度必须非负，并且不小于压强的大小。

\item \textbf{类光主能量条件}（Null Dominant Energy Condition），或称 NDEC，是仅对类光矢量成立的 DEC 条件：对任意类光矢量 $\ell^\mu$，$T_{\mu\nu}\ell^\mu\ell^\nu \geq 0$，并且 $T^\mu{}_\nu \ell^\nu$ 是非类空矢量。允许的密度和压强与 DEC 相同，只是允许负密度，只要 $p = -\rho$。换句话说，凡是会被 DEC 排除的源，NDEC 同样排除——负真空能除外。

\item \textbf{强能量条件}（Strong Energy Condition），或称 SEC，要求对所有类时矢量 $t^\mu$ 有 $T_{\mu\nu}t^\mu t^\nu \geq \frac{1}{2}T^\lambda{}_\lambda t^\sigma t_\sigma$，或者等价地，$\rho + p \geq 0$ 且 $\rho + 3p \geq 0$。注意 SEC 并\emph{不}蕴含 WEC。它蕴含 NEC，同时还排除了过大的负压强。从式 (4.86) 我们看到，正是 SEC 意味着引力是吸引的。
\end{itemize}

这些条件在图 4.3 中给出了解示。此外我们还画出了约束 $w \geq -1$，其中 $w = p/\rho$ 被称为\textbf{物态方程参数}（equation-of-state parameter）。这是宇宙学中一个有用的概念，因为宇宙学中的源常常具有 $p = w\rho$ 形式的物态方程，其中 $w$ 为常数（当然，无论 $w$ 是不是常数，它都是有定义的）。如果我们把自己限制在 $\rho \geq 0$ 的源，那么上述任何一个能量条件都将蕴含 $w \geq -1$。

%FIG{4.3}: {能量条件应用于理想流体的情形，表示为能量密度 $\rho$ 与压强 $p$ 平面上的允许区域。图中画出了弱能量条件（WEC）、类光能量条件（NEC）、主能量条件（DEC）、类光主能量条件（NDEC）以及强能量条件（SEC）。作为比较，我们还画出了条件 $w \geq -1$，其中 $w = p/\rho$ 为物态方程参数。}

大多数普通的经典物质形式，包括标量场和电磁场，都遵守 DEC（见习题），因而也遵守限制较弱的条件（WEC、NEC、NDEC）。SEC 在某些奇点定理的证明中很有用，但可以被某些形式的物质破坏，例如有质量的标量场。事实证明，量子场通常能够破坏我们列出的任何一种能量条件；不过，也可能存在一些涉及时空区域上积分的不等式，即使是量子场也满足它们。这是目前仍在活跃研究的一个领域。

严格说来，能量条件与能量守恒并没有关系；无论我们对 $T^{\mu\nu}$ 施加什么附加约束，Bianchi 恒等式都保证 $\nabla_\mu T^{\mu\nu} = 0$。它们的作用在于防止其他被我们视为“非物理”的性质，例如能量以超过光速的速度传播，或者空的空间自发衰变为正负能量互相补偿的区域。特别是，Hawking 和 Ellis (1973) 证明了一个\emph{守恒定理}（conservation theorem）：大致说来，如果能动张量遵守 DEC，并且在某个类空区域中为零，那么它必定在该区域的未来依赖域中处处为零（未来依赖域的定义见 2.7 节）。这样一来，能量不可能无中生有地冒出来，也不可能溜到光锥之外。该定理并不包括其逆命题（即违反 DEC 的源必然是非因果的），所以谨慎一点总是有好处的。

\section{再论等效原理}

本节我们将更仔细地考察等效原理的依据和后果，我们在 4.1 节中曾用它来启发把物理学推广到弯曲时空的最小耦合步骤。我们将看到，等效原理并不是一条神圣的物理定律，甚至算不上一个数学上严格的陈述；在更基本的层次上，它的产生是广义相对论作为一种在宏观距离上有效的有效场论这一本质的结果，而我们的任务则是确定在这样的理论中，我们应当预期物质与度规之间存在哪些类型的耦合。

在实践中，人们常常援引等效原理来为下面四种观念中的任何一种辩护：
\begin{enumerate}
\item 物理学定律应当（或至少可以）以广义协变的形式表述。
\item 时空中存在一个度规，其曲率被我们解释为引力。
\item 不存在任何其他与引力相像的场。
\item 物质场与曲率的相互作用是最小的：它们不包含与 Riemann 张量或其缩并的直接耦合。
\end{enumerate}

这些截然不同的陈述，各自有着非常不同的地位：第一种是空洞的，第二种既深刻又几乎肯定为真，第三种有趣且可检验，而第四种只不过是一个有用的近似。让我们逐一考察。

第一种陈述有时被称为\textbf{协变性原理}（Principle of Covariance）。它多少是没什么内容的。“广义协变”只是说，方程中所有的项在坐标变换下以相同的方式变换，从而方程的形式是坐标不变的。由于张量变换规律的普适性，实现这一目标最直接的办法，就是把方程写成明显的张量形式。当然，一条定律若以并非广义协变的形式表述，也没有任何问题，只要我们知道能够把它改写成与坐标无关的形式。另一方面，只要定律从一开始就是良定义的，把它写成与坐标无关的形式总是\emph{可能的}。一个以某种方式行为的物理系统，并不知道你正在用哪个坐标系来描述它；因此，任何配得上“物理学定律”这一名称的东西（相对于该定律的某个特定表述而言）都必定与坐标无关。坚持显式的坐标不变性，对于定律如何向弯曲时空推广并没有说什么；正如我们已看到的，明显的张量方程不论几何如何都取相同的形式。

考虑平直时空中的麦克斯韦方程组，如同我们在第 1 章中写下的那样：
\begin{equation*}
\partial_\mu F^{\nu\mu} = J^\nu.
\tag{4.109}
\end{equation*}
右边是一个良定义的张量，而左边由于出现了偏导数则不是。这没有问题，因为我们知道这个方程只在 Minkowski 空间的惯性坐标中有效。用坐标不变的方式表达同一定律，则是
\begin{equation*}
\nabla_\mu F^{\nu\mu} = J^\nu.
\tag{4.110}
\end{equation*}
要断言这就是 Minkowski 空间中正确的表述，无需援引任何物理原理；它是\emph{唯一的}张量方程，在惯性坐标中与式 (4.109) 等价。但它并不是向弯曲时空的唯一推广，因为我们可以想象一些新的项，涉及 $F_{\mu\nu}$ 与 $R^\rho{}_{\sigma\mu\nu}$ 的乘积；这类附加项的地位由最小耦合假设直接处理，即上面清单中的第四点。不过，“把东西弄成张量的”或“广义协变的”这件事本身，只是逻辑必然性的简单体现，并不是一个可以设想用实验来否证的物理原理。（同一思想的另一种说法是“微分同胚不变性”（diffeomorphism invariance），在附录 B 中讨论。）

上面清单中，等效原理的第二条所谓后果要深刻得多，而且绝非显然。虽然 Einstein 是受 EP 的启发，但这一几何洞见正是他的伟大突破。在第 2 章开头我们讨论过为什么这样的洞见是有根据的：EP 意味着引力是普适的，这又进而意味着引力场在小的时空区域内变得无法测量，而这一特征最直接的实现方式，就是把引力与时空几何的效应等同起来。这些步骤是有充分动机的建议，而不是严格推导出的结论；一旦我们有了“存在一个度规、其曲率产生引力”的想法，就可以通过与实验比较来检验它的用处。正如我们提到的，它以优异的成绩通过了检验。不断积累的证据（例如第 2 章讨论过的引力红移）与“理想化的直尺和时钟，其表现正如时空几何弯曲时应有的表现”这一想法相一致。尽管如此，不应指望去证明真的存在一个具有所期望性质的度规；我们提出假设，用越来越精确的实验去检验它，并推断它的适用范围。的确，最终调和广义相对论与量子力学的要求，让许多人相信度规终将被揭示为一个由更基本的一组自由度衍生出来的概念。就我们当前的目的而言，这个终极答案并不重要；弯曲度规这一想法的有用性已经被证明得无可置疑，我们要做的，是扩展对其性质的理解，直到碰上不可逾越的障碍（无论是理论上的还是经验上的）。

既然我们确信引力的效应最好归因于时空度规的曲率，那么如果实验探测到对等效原理的明显违反，我们会得出什么结论？例如，我们可能设想这样一个实验，它揭示检验物体朝向地球或太阳的加速度确实极其轻微地依赖于检验物体的成分。（目前对这类反常加速度最好的限制表明，它们小于引力所致加速度的 $10^{-12}$ 倍。）\footnote{Y.~Su et al., \emph{Phys. Rev.} \textbf{D~50}, 3614 (1994).} 在这种情况下，没有人会真的倾向于宣称广义相对论已被彻底推翻、有必要从头再来。更可能的做法是回到“检验物体”（test body）的定义上去，这个定义包含该物体不带电荷这一条件。举例来说，电子就不是一个好的检验物体，因为它除了受引力作用之外，还会被周围环境中的电磁场推来搡去。类似地，对于号称中性的检验物体上任何假想的反常加速度，迄今为止最直接的解释莫过于设想我们发现了一种新的长程场的存在，而我们的检验物体在这种场下其实是带荷的。这样的场要一直未被发现至今，它必须要么耦合极弱，要么必须近乎普适地耦合，以便模仿引力的效应。我们可以设想，比如，与能动张量的迹耦合的标量场，或者与重子数（baryon number）耦合的矢量场。普通检验物体的质量近似正比于其重子数，重子数计数的正是物体中质子和中子的数目。因此，有时把“等效原理的检验”方便地当作对我们上述第三条陈述的检验——即不存在任何其他与引力相像的场（如果一个场是长程的，并且近乎普适地与质量耦合，它就与引力相像）。同样，探测到对这一假设的违反，最直接的解释将是发现了一种新的“第五种力”（fifth force），而不是对 Einstein 思想的否定。至于如果我们改进现有实验，是否应当预期会发现这样一种新场，这很难说；一方面，随手就能构造出带有新长程力的模型，但另一方面，它们通常应当强到早已被探测到了。在这个阶段，保持开放的头脑仍然是值得的。

除了度规的存在本身之外，等效原理的核心在于我们表述中的第四条：物质场与曲率的相互作用是最小的，不包含与 Riemann 张量或其缩并的直接耦合。例如，我们可以考虑如下这个取代常规测地线方程的可能替代：
\begin{equation*}
\frac{d^2x^\mu}{d\lambda^2} + \Gamma^\mu_{\rho\sigma}\frac{dx^\rho}{d\lambda}\frac{dx^\sigma}{d\lambda} = \alpha\left(\nabla_\sigma R\right)\frac{dx^\mu}{d\lambda}\frac{dx^\sigma}{d\lambda},
\tag{4.111}
\end{equation*}
其中 $R$ 是 Ricci 标量曲率，$\alpha$ 是一个耦合常数。这个方程在平直时空中同样约化为直线运动，但通过测量与 $\nabla_\sigma R$ 的耦合，它将允许在小区域内直接探测时空曲率。那么，为什么大自然选择了简单的测地线方程呢？作为回答的第一步，考虑耦合 $\alpha$ 的量纲。由于 $c = 1$ 且空间和时间有相同的单位，我们可以用长度作为基本量纲。于是度规、逆度规以及 $dx^\mu/d\lambda$ 都是无量纲的。偏导数算符的单位是逆长度，协变导数也一样。Christoffel 符号包含度规的一阶导数，因此具有逆长度的量纲；类似地，Riemann 张量、Ricci 张量和 Ricci 标量曲率具有逆长度平方的量纲：
\begin{equation*}
\Big[\frac{dx^\mu}{d\lambda}\Big] = [g_{\mu\nu}] = [g^{\mu\nu}] = L^0, \qquad [\nabla_\mu] = [\Gamma^\mu_{\rho\sigma}] = L^{-1}, \qquad [R] = L^{-2}.
\tag{4.112}
\end{equation*}
为保持一致，耦合 $\alpha$ 必须具有长度平方的量纲：
\begin{equation*}
[\alpha] = L^2.
\tag{4.113}
\end{equation*}
因此 $\alpha$ 的平方根定义了一个长度尺度；这个长度尺度应当是多大？我们无法确定，但有充分的理由相信它应当极其小。有两点论证。其一，由于 $\alpha$ 所代表的耦合起源于引力，对相关长度尺度唯一合理的预期是
\begin{equation*}
\alpha \sim l_{\rm P}^2,
\tag{4.114}
\end{equation*}
其中 $l_{\rm P}$ 是 Planck 长度。另一个理由不过是“不然还能是什么呢”这种理由的更精致版本。虽然广义相对论是一个经典理论，但在更深的层次上，我们预期它只不过是描述某个底层量子力学结构的有效场论。即使不知道这个结构可能是什么，一个一般的预期（源自我们对已经理解的量子场论的经验）是：有效的经典极限应当包含所有可能的相互作用，只是带有量纲的长度参数，它们代表着新自由度开始变得重要的尺度（回忆第 1 章末尾我们对有效场论的讨论）。这样，弱相互作用的 Fermi 理论包含一个长度尺度，我们现在知道它对应于电弱对称性破缺的尺度，W 和 Z 玻色子在那里开始变得相关。既然我们不预期新的引力物理在 Planck 标度之前出现，与引力相伴的高阶相互作用就应当被 Planck 长度的适当幂次压低。

这代表着多大的压低？一种量度是将 $l_{\rm P}$（从而参数 $\alpha$ 可能的取值）与地球附近典型的引力长度尺度相比较。地球上引力的强度由重力加速度表征，$a_g = 980\ \mathrm{cm/sec^2}$。为构造一个具有长度量纲的量，我们定义
\begin{equation*}
l_\oplus = c^2/a_g \sim 10^{18}\ \mathrm{cm},
\tag{4.115}
\end{equation*}
其中符号 $\oplus$ 在这里代表地球（而不是直和）。于是，高阶引力效应的相对强度由下式量度：
\begin{equation*}
\frac{l_{\rm P}}{l_\oplus} \sim 10^{-51}.
\tag{4.116}
\end{equation*}
实际上，由于我们预期 $\alpha \sim l_{\rm P}^2$，压低将是 $10^{-102}$ 的量级。因此，似乎没什么必要担心这类耦合可能起的作用。但剧烈的偏离应当记在心里；近来关于大额外维的想法已经开启了在粒子加速器上观测直接引力相互作用的可能性。归根结底，没有办法仅仅靠纯粹的思考来解决这些问题；只有实验才能在各种替代方案中作出裁决。

\section{替代理论}

广义相对论已经通过了各种各样的实验检验。尽管如此，我们做的下一个实验仍有可能揭示对 Einstein 原始表述的偏离。因此，让我们简要考虑广义相对论可能被修改的几种方式。这样的方式多得不可胜数，但我们将考虑四种不同的可能性：
\begin{itemize}
\item 引力标量场
\item 额外空间维度
\item 作用量中的高阶项
\item 非 Christoffel 联络
\end{itemize}

一类流行的替代模型被称为引力\textbf{标量-张量理论}（scalar-tensor theories），因为它们同时涉及度规张量 $g_{\mu\nu}$ 和一个标量场 $\lambda$。特别地，标量场直接与标量曲率耦合，而不只是与度规耦合（正如等效原理似乎所蕴含的那样）。作用量可以写成引力部分、纯标量部分和物质部分之和：
\begin{equation*}
S = S_{fR} + S_\lambda + S_{\rm M},
\tag{4.117}
\end{equation*}
其中
\begin{equation*}
S_{fR} = \int d^4x\,\sqrt{-g}\, f(\lambda)R,
\tag{4.118}
\end{equation*}
\begin{equation*}
S_\lambda = \int d^4x\,\sqrt{-g}\left[-\frac{1}{2}h(\lambda)g^{\mu\nu}(\partial_\mu\lambda)(\partial_\nu\lambda) - U(\lambda)\right],
\tag{4.119}
\end{equation*}
以及
\begin{equation*}
S_{\rm M} = \int d^4x\,\sqrt{-g}\,\widehat{\mathcal{L}}_{\rm M}(g_{\mu\nu}, \psi_i).
\tag{4.120}
\end{equation*}
这里，$f(\lambda)$、$h(\lambda)$ 和 $U(\lambda)$ 是定义该理论的函数，而物质拉格朗日量 $\widehat{\mathcal{L}}_{\rm M}$ 依赖于度规和一组物质场 $\psi_i$，但不依赖于 $\lambda$。通过变量的改变，我们总可以令 $h(\lambda) = 1$，但我们把它保留在此，以便与文献中的模型进行比较。

这个理论的运动方程包括引力方程（由对度规变分得到）、标量方程（由对 $\lambda$ 变分得到），以及相应的物质方程。让我们从引力方程开始，我们可以按照与普通 Hilbert 作用量 (4.55) 相同的步骤来推导它。考虑度规的微扰
\begin{equation*}
g^{\mu\nu} \to g^{\mu\nu} + \delta g^{\mu\nu}.
\tag{4.121}
\end{equation*}
按照 4.3 节的步骤，作用量引力部分的变分为
\begin{align*}
\delta S_{fR} = {}& \int d^4x\,\sqrt{-g}\, f(\lambda)\Big[\Big(R_{\mu\nu} - \frac{1}{2}Rg_{\mu\nu}\Big)\delta g^{\mu\nu} + \nabla_\sigma\nabla^\sigma(g_{\mu\nu}\delta g^{\mu\nu}) \\
&- \nabla_\mu\nabla_\nu(\delta g^{\mu\nu})\Big].
\tag{4.122}
\end{align*}
对于 Hilbert 作用量，$f$ 是常数，所以最后两项是全导数，可以通过分部积分化为表面项从而忽略。现在，分部积分（两次）会得到 $f$ 的导数，我们便得出
\begin{equation*}
\delta S_{fR} = \int d^4x\,\sqrt{-g}\,\big[f(\lambda)G_{\mu\nu} + g_{\mu\nu}\Box f - \nabla_\mu\nabla_\nu f\big]\delta g^{\mu\nu},
\tag{4.123}
\end{equation*}
其中 $G_{\mu\nu}$ 是 Einstein 张量。我们像往常一样丢弃了表面项，尽管在这种情形下关于边界贡献有一些微妙之处；讨论见 Wald (1984)。于是，包含来自 $S_\lambda$ 和 $S_{\rm M}$ 的贡献的引力运动方程为
\begin{equation*}
G_{\mu\nu} = f^{-1}(\lambda)\left(\frac{1}{2}T^{({\rm M})}_{\mu\nu} + \frac{1}{2}T^{(\lambda)}_{\mu\nu} + \nabla_\mu\nabla_\nu f - g_{\mu\nu}\Box f\right),
\tag{4.124}
\end{equation*}
其中能动张量为 $T^{(i)}_{\mu\nu} = -2(-g)^{-1/2}\delta S_i/\delta g^{\mu\nu}$；特别地，
\begin{equation*}
T^{(\lambda)}_{\mu\nu} = h(\lambda)\nabla_\mu\lambda\nabla_\nu\lambda - g_{\mu\nu}\left[\frac{1}{2}h(\lambda)g^{\rho\sigma}\nabla_\rho\lambda\nabla_\sigma\lambda + U(\lambda)\right].
\tag{4.125}
\end{equation*}
观察式 (4.124) 中 $T^{({\rm M})}_{\mu\nu}$ 的系数可以看到，当标量场\emph{恒定}（或实际上如此）时，我们可以把 $f(\lambda)$ 认作 $1/(16\pi G)$，这从原始作用量 (4.118) 看是合理的。另一方面，如果 $\lambda$ 在时空中逐点略有变化，它将被解释为一个随时空变化的 Newton 引力常数。控制这种变化的动力学由 $\lambda$ 的运动方程确定，不难推导出为
\begin{equation*}
h\,\Box\lambda + \frac{1}{2}h'g^{\mu\nu}\nabla_\mu\lambda\nabla_\nu\lambda - U' + f'R = 0,
\tag{4.126}
\end{equation*}
其中撇号表示对 $\lambda$ 求导。注意，如果我们令 $h(\lambda) = 1$ 以得到标量的常规动能项，$\lambda$ 就服从常规的标量场运动方程，只是与标量曲率有一个附加耦合。在现实世界中，我们不希望 $f(\lambda)$ 变化太大，否则它将在太阳系中 GR 的经典实验检验、以及诸如原初核合成这样的宇宙学检验中产生可观测的后果。这可以通过两种方式得到保证：要么选择 $U(\lambda)$ 使势有一个极小值，$\lambda$ 若没有大的能量输入就不能偏离太远——换句话说，$\lambda$ 有很大的质量——要么选择 $f(\lambda)$ 和 $h(\lambda)$，使得 $\lambda$ 的大变化只引起 Newton 引力常数有效值的相对较小的变化。

最早的标量-张量模型之一被称为 Brans–Dicke 理论，在我们的记号下它对应于选取
\begin{equation*}
f(\lambda) = \frac{\lambda}{16\pi}, \qquad h(\lambda) = \frac{\omega}{8\pi\lambda}, \qquad U(\lambda) = 0.
\tag{4.127}
\end{equation*}
其中 $\omega$ 是一个耦合常数。标量-张量作用量取如下形式：
\begin{equation*}
S_{\rm BD} = \int d^4x\,\sqrt{-g}\left[\frac{\lambda}{16\pi}R - \frac{\omega}{16\pi}g^{\mu\nu}\frac{(\partial_\mu\lambda)(\partial_\nu\lambda)}{\lambda}\right].
\tag{4.128}
\end{equation*}
在 Brans–Dicke 理论中，标量场是无质量的，但在 $\omega \to \infty$ 的极限下，场变为非动力学的，通常的 GR 得到恢复。目前太阳系检验给出的界限意味着 $\omega > 500$，因此如果存在这样的标量场，它必定只能与 Ricci 标量弱耦合。

处理标量-张量理论的一种流行做法，是施行一个共形变换，把理论化成看起来像常规 GR 的形式。我们定义共形度规
\begin{equation*}
\widetilde{g}_{\mu\nu} = 16\pi\widetilde{G}f(\lambda)g_{\mu\nu},
\tag{4.129}
\end{equation*}
其中 $\widetilde{G}$ 将成为共形框架（conformal frame）中的 Newton 引力常数。利用附录 G 中的共形变换公式，来自 (4.118) 的作用量 $S_{fR}$ 变为
\begin{align*}
S_{fR} &= \int d^4x\,\sqrt{-g}\, f(\lambda)R \\
&= \int d^4x\,\sqrt{-\widetilde{g}}\,(16\pi\widetilde{G})^{-1}\left[\widetilde{R} - \frac{3}{2}\widetilde{g}^{\rho\sigma}f^{-2}\left(\frac{df}{d\lambda}\right)^2(\widetilde{\nabla}_\rho\lambda)(\widetilde{\nabla}_\sigma\lambda)\right],
\tag{4.130}
\end{align*}
其中像往常一样，我们已经分部积分并丢弃了表面项。于是，在共形框架中，标量曲率独自出现，不再乘以任何 $\lambda$ 的函数。这个框架有时被称为\textbf{Einstein 框架}（Einstein frame），因为 Einstein 方程对共形度规 $\widetilde{g}_{\mu\nu}$ 取其常规形式。带有度规 $g_{\mu\nu}$ 的原来框架被称为\textbf{Jordan 框架}（Jordan frame），有时也叫\textbf{弦框架}（string frame）。（弦论典型地预言的是标量-张量理论而非通常的 GR，而且弦的世界面对度规 $g_{\mu\nu}$ 有响应。）

在继续分析共形变换后的理论之前，考虑一下如果我们选取
\begin{equation*}
f(\lambda) = e^{\lambda/\sqrt{3}}, \qquad h(\lambda) = 0, \qquad U(\lambda) = 0,
\tag{4.131}
\end{equation*}
即 $f(\lambda)$ 的一个特定选择，但同时关掉 $S_\lambda$ 中的纯标量项，会发生什么。于是我们注意到，Einstein 框架的作用量 (4.130) 实际上包含了标量的常规动能项，尽管它并不出现在 Jordan 框架的作用量 (4.118) 中。即使没有 $\lambda$ 的显式动能项，这个理论的自由度也包含一个传播的标量以及度规。在我们于第 7 章考察引力场的自由度之后，这一点应该会变得更清楚。在那里我们将发现，度规 $g_{\mu\nu}$ 实际上包含标量（自旋 0）和矢量（自旋 1）自由度，以及预期的张量（自旋 2）自由度；然而，在标准 Hilbert 作用量下，这些自由度是受约束的，而不是自由传播的。我们刚刚发现的就是：在作用量中用一个标量去乘 $R$，其作用是让标量自由度活起来，这一点在 Einstein 框架中明白无误地显露出来。

如果我们确实选择把纯标量作用量 $S_\lambda$ 包括进来，就得到
\begin{equation*}
S_{fR} + S_\lambda = \int d^4x\,\sqrt{-\widetilde{g}}\left[\frac{\widetilde{R}}{16\pi\widetilde{G}} - \frac{1}{2}K(\lambda)\widetilde{g}^{\rho\sigma}(\widetilde{\nabla}_\rho\lambda)(\widetilde{\nabla}_\sigma\lambda) - \frac{U(\lambda)}{(16\pi\widetilde{G})^2 f^2(\lambda)}\right],
\tag{4.132}
\end{equation*}
其中
\begin{equation*}
K(\lambda) = \frac{1}{16\pi\widetilde{G}f^2}\big[fh + 3(f')^2\big].
\tag{4.133}
\end{equation*}

我们可以通过定义一个新的标量场 $\phi$，让我们的作用量看起来完全符合常规：
\begin{equation*}
\phi = \int K^{1/2}\,d\lambda,
\tag{4.134}
\end{equation*}
用它表出，作用量变为
\begin{equation*}
S_{fR} + S_\lambda = \int d^4x\,\sqrt{-\widetilde{g}}\left[\frac{\widetilde{R}}{16\pi\widetilde{G}} - \frac{1}{2}\widetilde{g}^{\rho\sigma}(\widetilde{\nabla}_\rho\phi)(\widetilde{\nabla}_\sigma\phi) - V(\phi)\right],
\tag{4.135}
\end{equation*}
其中
\begin{equation*}
V(\phi) = \frac{U(\lambda(\phi))}{(16\pi\widetilde{G})^2 f^2(\lambda(\phi))}.
\tag{4.136}
\end{equation*}
令人惊叹的是，在 Einstein 框架中我们得到的是一个完全平常的弯曲时空标量场理论。只要 $f(\lambda)$ 表现良好，变量 $(\widetilde{g}_{\mu\nu}, \phi)$ 就可以取代 $(g_{\mu\nu}, \lambda)$，其含义是：对新变量做变分，等价于从原来的运动方程 (4.124) 和 (4.126) 出发，然后施行变换 (4.129) 和 (4.134)。

最后，我们把物质作用量 (4.120) 加进来。对 $\widetilde{g}_{\mu\nu}$ 做变分将给出 Einstein 框架中的能动张量。在原变量 $(g_{\mu\nu}, \lambda)$ 下，我们知道 $S_{\rm M}$ 不依赖于 $\lambda$，但现在它将同时依赖于两个新变量 $(\widetilde{g}_{\mu\nu}, \phi)$；我们可以用链式法则来刻画这种依赖。让我们再假设 $S_{\rm M}$ 只是代数地依赖于 $g_{\mu\nu}$，而不通过其导数。对普通的标量场或规范场物质来说这是成立的；对费米子情形会变得更复杂，我们在此不予讨论。我们得到
\begin{align*}
\widetilde{T}_{\mu\nu} &\equiv -2\frac{1}{\sqrt{-\widetilde{g}}}\frac{\delta S_{\rm M}}{\delta\widetilde{g}^{\mu\nu}} \\
&= -2\frac{1}{\sqrt{-\widetilde{g}}}\frac{\partial g^{\alpha\beta}}{\partial\widetilde{g}^{\mu\nu}}\frac{\delta S_{\rm M}}{\delta g^{\alpha\beta}} \\
&= -2(16\pi\widetilde{G}f)^{-1}\frac{1}{\sqrt{-g}}\delta^\alpha_\mu\delta^\beta_\nu\frac{\delta S_{\rm M}}{\delta g^{\alpha\beta}} \\
&= (16\pi\widetilde{G}f)^{-1}T_{\mu\nu}.
\tag{4.137}
\end{align*}
同样的技巧也适用于物质与 $\phi$ 的耦合，它来自对 $\phi$ 变分 $S_{\rm M}$，利用 $g^{\alpha\beta} = 16\pi\widetilde{G}f\,\widetilde{g}^{\alpha\beta}$：
\begin{align*}
\frac{\delta S_{\rm M}}{\delta\phi} &= \frac{\partial g^{\alpha\beta}}{\partial\phi}\frac{\delta S_{\rm M}}{\delta g^{\alpha\beta}} \\
&= \left(16\pi\widetilde{G}\frac{df}{d\phi}\widetilde{g}^{\alpha\beta}\right)\left(-\frac{1}{2}\sqrt{-g}\,T^{\rm M}_{\alpha\beta}\right) \\
&= -\frac{1}{2f}\frac{df}{d\phi}\sqrt{-\widetilde{g}}\,\widetilde{T}^{\rm M},
\tag{4.138}
\end{align*}
其中
\begin{equation*}
\widetilde{T}^{({\rm M})} = \widetilde{g}^{\alpha\beta}\widetilde{T}^{({\rm M})}_{\alpha\beta} = \frac{1}{(16\pi\widetilde{G}f)^2}g^{\alpha\beta}T^{({\rm M})}_{\alpha\beta}
\tag{4.139}
\end{equation*}
是共形框架中能动张量的迹。

对 $\widetilde{g}_{\mu\nu}$ 和 $\phi$ 变分 (4.135)，得到与 Einstein 方程等价的运动方程以及一个 $\phi$ 的方程。引力方程是
\begin{equation*}
\widetilde{G}_{\mu\nu} = 8\pi\widetilde{G}\left(\widetilde{T}^{({\rm M})}_{\mu\nu} + \widetilde{T}^{(\phi)}_{\mu\nu}\right),
\tag{4.140}
\end{equation*}
其中
\begin{equation*}
\widetilde{T}^{(\phi)}_{\mu\nu} = \widetilde{\nabla}_\mu\phi\,\widetilde{\nabla}_\nu\phi - \widetilde{g}_{\mu\nu}\left[\frac{1}{2}\widetilde{g}^{\rho\sigma}\widetilde{\nabla}_\rho\phi\,\widetilde{\nabla}_\sigma\phi + V(\phi)\right],
\tag{4.141}
\end{equation*}
而标量场方程是
\begin{equation*}
\widetilde{\Box}\phi - \frac{dV}{d\phi} = \frac{1}{2f}\frac{df}{d\phi}\widetilde{T}^{({\rm M})}.
\tag{4.142}
\end{equation*}

既然 (4.140) 看起来就像同时带有物质源和标量场源的 Einstein 方程，为什么我们还要费心把这个标量-张量理论称作 GR 的一种替代理论呢？它难道不就是同一个理论，只是用了不同的变量吗？事实上它并不是同一个理论，原因在于 Einstein 框架中 $S_{\rm M}$ 对 $\phi$ 的依赖。特别是，物理的检验粒子将沿着 $g_{\mu\nu}$ 的测地线运动，而这些测地线一般不会与 $\widetilde{g}_{\mu\nu}$ 的测地线重合。检验粒子“看到”的是原来的度规。所以，要么我们在原变量 $(g_{\mu\nu}, \lambda)$ 中工作，此时引力场方程被修改；要么我们使用新变量 $(\widetilde{g}_{\mu\nu}, \phi)$，此时物质的运动方程被修改；无论哪种方式，都会存在（原则上）毫不含糊、可测量的对通常 GR 的偏离。

修改广义相对论的另一种途径，是允许额外空间维度的存在；事实上，额外维的物理后果与标量-张量理论的物理后果是紧密相关的。所谓额外维，我们并不是简单地在高维空间中考虑 GR，而是考虑这样一类模型：时空在大尺度上表现为四维，尽管实际上总共有 $4 + d$ 维。实现这一点最简单的方式，是让额外的 $d$ 维在某一流形上“紧化”（compactified）；我们在这里考虑的正是这种可能性。\footnote{我们沿用 S.M.~Carroll, J.~Geddes, M.~Hoffman, and R.M.~Wald 的分析，\emph{Phys. Rev.} \textbf{D~66}, 024036 (2002)；http://arxiv.org/hep-th/0110149。关于额外维的原始论文出自 Kaluza 和 Klein：T.~Kaluza, \emph{Sitzungsber. Preuss. Akad. Wiss. Berlin (Math. Phys.)} K1, 966 (1921)；O.~Klein, \emph{Z. Phys.} \textbf{37}, 895 (1926) [\emph{Surveys High Energ. Phys.} \textbf{5}, 241 (1926)]；O.~Klein, \emph{Nature} \textbf{118}, 516 (1926)。}这类模型被称为 Kaluza–Klein 理论。

令 $G_{ab}$ 为一个 $(4 + d)$ 维时空的度规，坐标为 $X^a$，指标 $a, b$ 从 0 取到 $d + 3$。
